\chapter{Conclusions and Future Work}
\label{ch:conclusions}
\section{Conclusions}

This thesis implements and evaluates LM-IDNet, a normal-only anomaly
detector for device-specific IoT traffic. Its input is a completed
window of TCP, ordinary UDP, SSDP, and ARP counts. A Dirichlet--multinomial
model supplies a probabilistic reference for composition; separately
calibrated count features complement that reference within one joint
decision rule. The work connects this statistical procedure to
validated capture preparation, native numerical computation, guarded
candidate updates, persistent state, and a supplied-window runtime.

The foundations are established methods. The DM distribution,
Minka fixed-point estimator, and Languasco--Migliardi numerical
method are attributed to their sources. The contribution is their
working integration and the evidence collected about its behaviour,
including negative results. It does not require claiming a new
distribution, an original estimator, or the first use of learned
normal IoT profiles.

The detection evaluation demonstrates both promise and limits.
In the fixed closing replay, 1,008 unique labelled windows yield
14 true positives, three false positives, four false negatives,
and 987 true negatives. Pooled precision is 82.35\%, recall 77.78\%,
and F1 80.00\%. Device results differ substantially: Chromecast
F1 is 61.54\% and Samsung F1 90.91\%. Eight of nine fully observed
attack intervals are detected, with seven receiving an alert by
attack end. These are descriptive results on an inspected temporal
partition, not an untouched external performance estimate.

The duration study confirms that observation granularity belongs
in the interpretation of a detector. Coarser windows can improve
their own F1 while increasing delay and reducing the number of
alerts available before attack end. A high coarse-window score
therefore cannot establish prompt response. The retained ten-minute
reference is evaluated and operationally explicit, without being
claimed as a universally optimal duration.

The numerical investigation preserves the distinction between
estimation and likelihood evaluation. LM and SciPy produce no
different development decisions in the three-dataset comparison,
and their fitted quantities agree closely. LM is faster for the
tested single-window operation, while SciPy is faster for the
production matrix kernel and full fitting. The aggregate experiment
has a different workload and favours LM there. This mixed result
supports choosing a backend using the delivered operation, rather
than assuming that one method is fastest throughout the application.

Adaptation is the least resolved statistical component. A historical
refit improves Chromecast development detection but fails to improve
Samsung and permits large concentration changes. Conservative guards
reject every closing replay candidate, leaving the three modes
equivalent to static operation. Other-duration threshold updates
can still increase false alerts sharply. The supported achievement
is explicit, inspectable control over candidate admission and
activation, rather than demonstrated robust learning under arbitrary
benign or adversarial drift.

The operational evaluation shows that the supplied-window process
fits inside a one-CPU, 512~MiB Docker allocation on the measured host.
The long synthetic replay scores 16,000 of 20,000 input records,
with mean processing near 30~ms and process high-water RSS near
69~MiB. Refit rejection contributes the largest processing delays,
up to about 5.08~s. These results support local feasibility for the
tested interface; they do not include a physical gateway or live
packet capture.

Reproducibility is part of this outcome. Configurations, prepared
inputs, model identities, saved events, candidate histories, and
source fingerprints connect the narrative to the experiments.
Generated tables recompute reported metrics from retained counts.
Later scoring changes have an explicit regression boundary, rather
than being silently folded into historical measurements. The final
software verification passes all 643 current tests, supporting the
implemented contracts without substituting for detection evidence.

\section{Principal limitations}

The main empirical limitation is the absence of a newly untouched
evaluation after all design choices were completed. Later labels
were inspected in earlier audits, and a fixed closing replay cannot
undo that exposure. There are only two labelled device views and
18 positive closing windows. The D-Link workload provides useful
numerical evidence but has no attack labels, and repeated fit-window
bounds limit its chronological interpretation.

The representation cannot distinguish packet sequences that share
the same four counts. Low-rate attacks, TCP flag changes without
count changes, and other semantic differences may remain invisible.
Normal-only learning also assumes an adequate benign reference;
neither the initial fit nor the runtime buffer proves that assumption.
Accepted adaptation is not independently validated on its own
learning data, and stepwise movement limits do not guarantee
long-term detection quality.

Operational evidence is limited to a finite synthetic supplied-window
simulation and software tests. Memory stability over that run is
not an indefinite guarantee. Collector correctness, physical storage
faults, embedded hardware performance, energy use, remote event
delivery, and automatic response are not evaluated. These limits
define what remains to establish before describing the prototype
as a complete deployed intrusion-detection service.

\section{Future work}
\label{sec:future-work}

\subsection{Independent evaluation before further tuning}

The first priority is an untouched temporal extension or external
collection evaluated after fixing the representation, calibration,
and update policy. Additional manufacturers, firmware versions,
benign uses, and attack families should be represented. Device and
capture blocks should be retained as evaluation units so that
variation is measured across independent situations, rather than
through repeated outputs on the same windows.

That study should keep window confusion counts, attack-interval
coverage, alert availability, and observed exposure together.
An external benchmark against other detectors would require a
shared preprocessing and tuning protocol. Correcting overlapping
D-Link fit bounds and repeating both numerical backends would also
provide a clean chronological workload, clearly identified as a
new experiment rather than a revision of old measurements.

\subsection{Controlled drift and contamination experiments}

Adaptation needs known benign changes and controlled attack
contamination to separate their effects. A useful experiment would
vary the contamination fraction and the rate of benign change,
then measure buffer membership, concentration, threshold trajectory,
candidate rejection, and detection against a static reference.
Independent validation observations for candidate promotion would
avoid treating increased likelihood on the fitting buffer as
evidence of improved future performance.

Possible policy changes include an explicit cumulative movement
budget relative to the original state and calibration-tail
requirements scaled to available sample size. Their value must
be tested, particularly against the observed Samsung threshold
failure, rather than assumed from stricter-looking equations.
The useful target is successful adaptation to a verified benign
change while retaining attack sensitivity, not a larger number
of accepted updates.

\subsection{Representation and response granularity}

Compact TCP flag, direction, endpoint, or packet-size summaries
could distinguish behaviours aliased by the four-category input.
Their inclusion would require renewed preparation from source
captures and explicit feature contracts. Ablation experiments
should establish whether each added signal improves detection
or delay enough to justify its collection and calibration cost.

Shorter or multiscale observations could reduce waiting for an
aggregate to close. They would also change sparsity, temporal
dependence, sample counts, and false-alert exposure. A comparison
must therefore retrain and recalibrate, keep common support,
and use actual alert-availability times. Merely projecting an
hourly decision backward cannot demonstrate earlier detection.

\subsection{Live and physical deployment}

A live collector should first demonstrate byte-to-count agreement
with the offline preparation on a shared replay. A subsequent
physical-gateway study can measure collection losses, total
resource use, energy, restart behaviour, and delay through the
complete observation-to-alert path. Multi-device operation and
upstream backpressure need explicit load tests rather than an
extrapolation from one detector process.

Finally, remote event retention and operator response need defined
delivery and acknowledgement contracts. The existing atomic state
mechanism provides a sound local boundary for accepted parameters,
but a usable monitoring service also needs trustworthy input and
durable handling of its output. Evaluating those boundaries would
turn the present reproducible prototype into evidence for a
specific operational deployment.
